\section{Proof of exact attainment}
\label{V:exact-proof}

The preceding construction depends only on the curvature-null field.
We now place the given divisors in a two-dimensional section ring and
apply Proposition~\ref{prop:geometric}.

\begin{proof}[Proof of Theorem~\ref{V:exact}]
Choose two input divisors with distinct exponents, write them as
$N_1,N_2$, and denote their exponents by $u_1,u_2$. For a rational
vector $t=(t^1,t^2)$ put $D(t)=t^1N_1+t^2N_2$. Let $t_j$ solve
\[
 u_1t_j^1+u_2t_j^2=1,\qquad t_j^1+t_j^2=1/m_j.
\]
After discarding finitely many terms, these points approach from one
side the rational point $t_*$ given by the same first equation and
$t_*^1+t_*^2=0$. Hence $D(t_*)\sim_{\Q}L$, and there are rational
powers of rational functions $h_j$ such that
\begin{equation}\label{V:principal-corrections}
 N_j/m_j=D(t_j)+\divisor_X(h_j).
\end{equation}
A positive power clearing denominators is an ordinary rational function.
Its polar divisor is bounded by a positive linear combination of $N_1,N_2$,
because the left side of \eqref{V:principal-corrections} is effective.
That bound is null. Thus the power belongs to $k$, and we may regard
$h_j$ as an element of $k^\times\otimes_{\Z}\Q$.

If $k=\C$, the principal correction vanishes. The fixed-support limit
of the effective divisors $D(t_j)$ is effective, and linearity on that
finite support gives
\[
 v(D(t_*))=\lim_j v(N_j)/m_j=0.
\]
Then $N=D(t_*)$ proves the assertion.

Otherwise use the equidimensional model of Section~\ref{V:null-field}.
Define the divisor $P(t)$ on $Y$ by
\[
 P(t)=\sum_Q\min_{E:\,f(E)=Q}
 \frac{\coeff_E p^*D(t)}{\ord_E f^*Q}\,Q.
\]
Its support is finite and fixed, since the support of $p^*D(t)$ lies
in that of $p^*N_1+p^*N_2$. Each coefficient is a minimum of finitely
many rational linear functions. On a sufficiently short rational
segment $[t_*,t_0]$ containing a tail of $t_j$, all these minima are
linear. Reparameterize this segment by $u\in[0,1]$, so that $u=0$
corresponds to $t_*$ and $u=1$ to $t_0$.

Apply the correction-ring theorem to $P(t_0),P(t_*)$. It gives a
finitely generated bigraded section ring for
\[
 B_1=k_0(A_0+\epsilon P(t_0)),\qquad
 B_0=k_0(A_0+\epsilon P(t_*)).
\]
All the hypotheses of Proposition~\ref{prop:geometric} now hold:
$p_*f^*A_0=0$ by the toroidal-volume theorem; the normalized effective
divisors in \eqref{V:principal-corrections} approach $D(t_*)$ through
principal corrections from $Y$; their valuations tend to zero; and
the minimum divisor is affine on the chosen segment. That proposition
gives an effective rational divisor $N\sim_{\Q}D(t_*)\sim_{\Q}L$
with $v(N)=0$.
\end{proof}

In particular, no valuation of the exceptional correction is assumed to
vanish. The proof of Proposition~\ref{prop:geometric} first removes it
by pushforward and then uses centeredness on $X$. This is also why the
fixed valuation may have center of any codimension.
